% Part: lambda-calculus
% Chapter: introduction
% Section: lambda-computable

\documentclass[../../../include/open-logic-section]{subfiles}

\begin{document}

\olfileid{lam}{int}{cmp}
\olsection{\usetoken{S}{lambda definable} फलने संगणनक्षम असतात}

\begin{thm}
\ollabel{thm:lambda-computable}
आंशिक फलन $f$ हे एखाद्या लॅम्डा पदाने !!{lambda defined}
असेल, तर ते संगणनक्षम आहे.
\end{thm}

\begin{proof}
फलन~$f$ हे लॅम्डा पद~$X$ ने !!{lambda defined} आहे असे समजा.
~$f$ चे संगणन करण्याची एक अनौपचारिक पद्धत सांगू.
$m_0$, \dots,~$m_{n-1}$ ही निविष्ट मिळाल्यावर
$X \num m_0 \ldots \num m_{n-1}$ हे पद लिहा आणि आधी तेच चर्च अंक आहे का ते तपासा.
मग मूळ पदाची सर्व एका-पायरीतील न्यूनीकरणे लिहून वृक्ष तयार करा;
त्याखाली त्या प्रत्येक पदाची सर्व एका-पायरीतील न्यूनीकरणे
(म्हणजे मूळ पदाची दोन-पायरीतील न्यूनीकरणे) लिहा; आणि
प्रत्येक पातळी पूर्ण करून असेच पुढे चालू ठेवा. पदातील redex
जागा सांत असतात; प्रत्येक पायरीतील आवश्यक नामांतरासाठी
ठरलेली ताजी नावे घेतल्याने प्रत्येक पातळीही सांत राहते.
नामांतरापर्यंत कोणताही चर्च अंक मिळाला, तर संबंधित संख्या
उत्तर म्हणून परत करा. फलनाचे मूल्य अपरिभाषित असेल, तर
असा अंक मिळत नाही आणि शोध अखंड चालू राहतो.

चर्च यांच्या प्रबंधाचा आधार घेतल्यास हे फलन संगणनक्षम
असल्याचे कळते. प्रमेय सिद्ध करण्याचा अधिक चांगला मार्ग
म्हणजे या शोधपद्धतीचे पुनरावर्ती वर्णन देणे. उदाहरणार्थ,
“$\fn{IsASubterm}$,” “$\fn{Substitute}$,”
“$\fn{ReducesToInOneStep}$,” “$\fn{ReductionSequence}$,”
“$\fn{Numeral}$” इत्यादी आदिम पुनरावर्ती फलने आणि
संबंधांची मालिका परिभाषित करता येईल. मग
$f(m_0, \dots, m_{n-1})$ चे संगणन करणारी आंशिक
पुनरावर्ती पद्धत म्हणजे $X \num{m_0} \dots \num{m_{n-1}}$
पासून सुरू होऊन अंकापाशी संपणारा एका-पायरीतील
न्यूनीकरणांचा अनुक्रम शोधणे आणि त्या अंकाशी संबंधित
संख्या परत करणे. तपशील मोठे आणि कंटाळवाणे आहेत,
पण इतरथा नेहमीचे आहेत.
\end{proof}

\end{document}
